sexualiteit en spirualiteit
Is Seks Verenigbaar met een Spiritueel Leven?
Is seks onverenigbaar met een religieus leven? Welke plaats heeft de menselijke relatie in spirituele inspanningen?
Laten we allereerst de vraag stellen: waarom hebben mensen over de hele wereld seks zo’n enorm belangrijke plaats in hun leven gegeven? Waarom is dit zo? Zeker nu, in de huidige tolerante samenleving, waar jongens en meisjes van twaalf of dertien al seksueel actief zijn. Waarom hebben mensen, door al hun bezigheden en gedurende hun hele leven, dit ene aspect zo’n kolossale importantie gegeven?
Voordat we antwoorden, moeten we deze vraag samen onderzoeken. Deel de vraag met mij. Je luistert niet zomaar naar een of ander orakel, maar samen gaan we op onderzoek uit. Het is jouw leven waar we naar kijken.
Seks en Religie: Twee Uitersten
Er zijn goeroes en een hele filosofie genaamd Tantra, die deels gebaseerd is op seks. Volgens deze leer kun je via seks God bereiken, wat die God dan ook moge zijn. Dit is een zeer populaire opvatting.
Aan de andere kant heb je degenen die seks volledig hebben afgezworen, zoals de monniken, de Indiase sannyasins en de boeddhistische priesters. Zij hebben altijd volgehouden dat seks een verspilling van energie is. Om God te dienen, zo geloven zij, moet je al je energie bundelen. Daarom moet je seks ontkennen, onderdrukken en de innerlijke verlangens die in je branden de baas zijn. Controleer het.
Zo heb je dus twee uitersten: de permissieve, tolerante benadering en de zogenaamde religieuze onderdrukking. En daartussenin bevinden zich de mensen die van alles proberen te genieten. Ze hebben één voet in de ene wereld en de andere voet in de andere. Ze proberen over beide te praten en te zien of ze de twee met elkaar kunnen harmoniseren om zo God, of wat dan ook, te vinden. Uiteindelijk zullen ze waarschijnlijk vooral veel problemen vinden.
Waarom die Enorme Nadruk op Seks?
Dus, we vragen ons af: waarom hebben man en vrouw van deze hele ‘seks-business’ zo'n belangrijk punt gemaakt? Waarom geef je niet dezelfde importantie aan liefde, aan compassie, aan het principe om niet te doden? Waarom hecht je zo’n immense waarde uitsluitend aan seks?
Volg je wat ik zeg? Denk aan jullie oorlogen, terreur, nationale verdeeldheid, de hele immorele samenleving waarin je leeft. Waarom geef je niet evenveel belang aan al die zaken, in plaats van alleen aan dit ene ding? Volg je mijn vraag?
De Grootste Vreugde in een Zinloos Bestaan
Is het misschien omdat seks je grootste genot in het leven is? De rest van je leven is een sleur, een zwoegen, een strijd, een conflict, een betekenisloos bestaan. En dit ene ding geeft je tenminste een zeker gevoel van groot genot, een gevoel van welzijn, een gevoel van wat je relatie en ook wat je liefde noemt. Is dat de reden waarom we zo seksueel bezeten zijn?
Geef antwoord aan jezelf. Het is omdat we op geen enkel ander vlak vrij zijn. We moeten naar kantoor van negen tot vijf, waar je wordt gepest, waar bazen over je heersen. Je kent het wel, alles wat er op een kantoor, in een fabriek of bij een andere baan gebeurt, waar altijd iemand is die je domineert.
Onze geest is daardoor mechanisch geworden. We herhalen, herhalen, herhalen. We vallen in een traditie, in een sleur, in een routine. Ons denken is: "Ik ben een christen, ik ben een boeddhist, ik ben een hindoe, ik ben een katholiek, ik aanbid de paus." Het hele pad is duidelijk uitgestippeld en we volgen het. Of je verwerpt dat allemaal en vormt je eigen routine. Onze geest is dus een slaaf geworden van verschillende bestaanspatronen. Het is mechanisch geworden.
En seks, hoewel plezierig, wordt daardoor ook mechanisch.
Is Liefde Seks, Genot of Herinnering?
Als je hier dieper op in wilt gaan, moet je jezelf de vraag stellen: is liefde seks? Vraag het jezelf af. Is liefde genot? Is liefde verlangen? Is liefde de herinnering aan een gebeurtenis die je seks noemt, met alle verbeelding, de beelden, het erover nadenken? Is dat liefde? In godsnaam, is liefde een herinnering?
De oorspronkelijke vraagsteller vroeg ook welke plaats de menselijke relatie heeft in spirituele inspanningen. Kijk waar de menselijke relatie toe is gereduceerd: genot, seks, conflict, ruzies, verdeeldheid. Jij gaat jouw weg, ik ga mijn weg. Dat is onze feitelijke relatie in ons dagelijks leven. En welke plaats heeft zo’n relatie in een spirituele zoektocht?
Het is duidelijk dat de huidige vorm van relaties daar absoluut geen plaats in heeft. We zijn jaloers op elkaar, we willen elkaar bezitten, we willen elkaar domineren. Daardoor is er vijandigheid tussen ons. Als we seksueel onbevredigd zijn, zoeken we het elders. En zelfs binnen die seksuele relatie is er eenzaamheid, en het voortdurend zoeken naar je eigen genot.
Is dat allemaal liefde? We negeren datgene wat we liefde noemen, terwijl dat misschien wel het meest wonderbaarlijke is wat iemand kan hebben. We zijn zo verstrikt geraakt in de draaikolk van ons eigen genot. We willen niet alleen seksuele bevrediging, maar bevrediging op elk gebied. En dat noemen we liefde. Voor die ‘liefde’ zullen we elkaar zelfs doden, zoals bij de liefde voor het vaderland.
Het Verlies van Creatief Vermogen
Denk hier eens over na als dit afgelopen is. Waarom hebben man en vrouw dit ene ding zo buitengewoon belangrijk gemaakt? Alle tijdschriften, je weet wel wat er allemaal gebeurt. Waarom?
Is het omdat man en vrouw hun creatieve vermogen hebben verloren? Ik heb het niet over seksueel vermogen, maar over creatief vermogen. Het vermogen om te zien, om een licht voor zichzelf te zijn, om niemand te volgen, om geen enkel beeld, illusie of geloof te aanbidden.
Wanneer je dat allemaal opzij zet, en je je eigen kleine verlangens – je seksuele eisen en bevredigingen – hebt overstegen, wanneer je dat allemaal inziet, dan ontstaat er creatie. Dat betekent niet een schilderij maken of een gedicht schrijven. Het is dat gevoel van frisheid, van een geest die continu fris, jong en onschuldig is, niet vertroebeld en belast door allerlei herinneringen, ontevredenheden, angsten en zorgen.
Wanneer je dat allemaal kwijt bent, ontstaat er een totaal ander soort geest. Dan heeft seks zijn eigen, juiste plaats. Maar wanneer seks een middel wordt voor spirituele inspanningen, dan raken we compleet vast.
De Noodzaak van Twijfel
Blijkbaar missen we de kwaliteit van scepsis. Het vermogen om sceptisch te zijn over onze eigen verlangens, om de talloze goeroes te bevragen en in twijfel te trekken. Maar twijfel wordt ook gevaarlijk als je het niet in toom houdt. Dan ga je aan alles twijfelen en is er geen einde aan. Het is als een hond aan de lijn; je moet hem af en toe loslaten, zodat hij kan genieten en rondrennen. Op dezelfde manier moet twijfel aan een lijn worden gehouden, maar je moet de lijn ook af en toe loslaten. Zo blijft de geest – je hart, je brein, je emoties, alles – actief, en niet alleen maar gericht op die ene richting: seks, seks, seks.
